\documentclass[11pt]{article}
\usepackage{amsmath,amssymb}
\newcommand{\D}[2]{|D_{#1}^{(#2)}\rangle}
\begin{document}
\title{Exact bridges between optimal AC subspaces and highly AC pure states}
\author{Catalogue calculation}\date{31 October 2026}\maketitle

These identities relate existing exact representatives; they do not assert a
classification or uniqueness of the optimal subspaces. All coefficients are in
the normalized Dicke basis. The accompanying \texttt{bridge-verify.py} checks
the formulas with exact symbolic arithmetic against the saved catalogue bases.

\paragraph{A centered-projector construction.}
Let $P$ be a rank-$k$ orthogonal projector on the spin-$n/2$ space, of
dimension $d=n+1$, with $0<k<d$, and suppose that
\[
Y=P-\frac{k}{d}I_d
\]
has only irreducible tensor rank $n$. Set
$F_{rs}=(-1)^r\delta_{s,n-r}$ and use the row-column vectorization convention
$\operatorname{vec}(A)=\sum_{r,s}A_{rs}\D{n}{r}\D{n}{s}$. Then
\[
|\Phi_P\rangle=\frac{\operatorname{vec}(YF)}{\sqrt{\operatorname{Tr}Y^2}}
\]
is a symmetric $2n$-qubit pure state of AC order at least $n-1$.

Indeed, multiplication by $F$ intertwines conjugation of operators with the
tensor-product spin action, so tensor rank $n$ becomes total spin $n$, the
fully symmetric $2n$-qubit sector. Projector idempotency gives
\[
Y^2=\left(1-\frac{2k}{d}\right)Y+\frac{k(d-k)}{d^2}I_d,
\qquad \operatorname{Tr}Y^2=\frac{k(d-k)}d.
\]
Consequently its $n$-qubit reduced state is
\[
\rho_n=\frac{Y^2}{\operatorname{Tr}Y^2}
=\frac{I_d}{d}+\frac{d-2k}{k(d-k)}Y.
\]
Partial trace onto $n-1$ symmetric qubits is rotation covariant, and its
output operator space has no rank-$n$ tensors. It therefore kills $Y$ and
gives $\rho_{n-1}=I_n/n$, proving the claim. This is a sufficient
construction condition on a projector, not a condition satisfied by every
AC subspace.

This sufficient condition is impossible for $n\ge7$. Indeed, its output
would have $2n$ qubits and AC order $n-1$. If $n=2r+1\ge7$, the pure
even-order bound requires $2n=4r+2\ge(r+1)^2$, false for $r\ge3$.
If $n=2r+2\ge8$, the odd-order bound requires
$2n=4r+4\ge(r+1)(r+2)$, likewise false for $r\ge3$.
Thus the following bridges are exceptional finite constructions, not the
start of an unrestricted family under this same sufficient condition.

\paragraph{Four qubits to eight.}
The catalogue's optimal two-dimensional 1-AC subspace on four qubits has
orthonormal basis
\[
u_\pm=\frac12\D4{0}\pm\frac{i}{\sqrt2}\D4{2}+\frac12\D4{4}.
\]
For $P_4=|u_+\rangle\langle u_+|+|u_-\rangle\langle u_-|$, the operator
$Y_4=P_4-2I_5/5$ has only rank four. The construction gives
\[
\boxed{|\Phi_8\rangle=
\frac{\sqrt{30}}{12}\bigl(\D8{0}+\D8{8}\bigr)
+\frac{\sqrt{21}}6\D8{4}.}
\]
Its reduction is exactly $\rho_3=I_4/4$. The Majorana polynomial is
proportional to $z^8+14z^4+1$, the standard cubic constellation. Here the
polynomial convention is $\sum_r(-1)^r\sqrt{\binom Nr}c_rz^r$.

\paragraph{The six-, seven-, and twelve-qubit bridge.}
The published two-dimensional 2-AC subspace on seven qubits (Eq.~E1 of
PRA 111, 022435 (2025), catalogue \texttt{paper-e1}) is spanned by
\[
u=\sqrt{\frac3{10}}\D7{0}+\sqrt{\frac7{10}}\D7{5},\qquad
v=\sqrt{\frac7{10}}\D7{2}-\sqrt{\frac3{10}}\D7{7}.
\]
Write $P_7=|u\rangle\langle u|+|v\rangle\langle v|$ and let $\mathcal R_6$
trace out one qubit. Its Dicke-basis matrix is
\[
(\mathcal R_6 X)_{rs}=
\frac{\sqrt{(7-r)(7-s)}X_{rs}
+\sqrt{(r+1)(s+1)}X_{r+1,s+1}}7.
\]
Direct substitution proves that $Q=2\mathcal R_6(P_7)$ is an orthogonal
rank-four projector. More strongly,
\[
\boxed{P_6=I_7-2\mathcal R_6(P_7),}
\]
where $P_6$ projects onto the catalogue's optimal three-dimensional 1-AC
subspace with orthonormal basis
\[
w_1=\sqrt{\frac35}\D6{1}+\sqrt{\frac25}\D6{6},\quad
w_2=\sqrt{\frac25}\D6{0}-\sqrt{\frac35}\D6{5},\quad
w_3=\D6{3}.
\]
Thus the \emph{kernel} of the reduced normalized code state
$\mathcal R_6(P_7/2)=Q/4$ is an optimal 1-AC subspace. No such conclusion
is claimed for its support.

The centered projectors obey
\[
Y_6=P_6-\frac37I_7=-\left(Q-\frac47I_7\right),
\qquad Y_6^2=\frac17Y_6+\frac{12}{49}I_7.
\]
$Y_6$ has only tensor rank six. Vectorizing it as above gives, up to a
global sign,
\[
\boxed{|\Phi_{12}\rangle=
\frac{-\sqrt7\D{12}{1}+\sqrt{11}\D{12}{6}
+\sqrt7\D{12}{11}}5.}
\]
Its reduction is exactly $\rho_5=I_6/6$. Therefore the three exact catalogue
maxima $(N,t,k)=(6,1,3),(7,2,2),(12,5,1)$ share an explicit algebraic
construction. The Majorana polynomial is proportional to
$z(1+11z^5-z^{10})$, with an additional point at infinity, a standard
icosahedral form.

\paragraph{Exact checks.}
Besides every ordered cross reduced matrix of the three source subspaces,
the verifier checks projector idempotency, the complement identity, the
rank-four/rank-six tensor claims via
$\sum_a[J_a,[J_a,Y]]=L(L+1)Y$, and every entry of both output states'
reduced matrices. It also checks directly that the vectorized matrix lies
in the highest-spin sector using the Dicke Clebsch--Gordan coefficients
$\sqrt{\binom nr\binom ns/\binom{2n}{r+s}}$.
\end{document}
